\chapter*{Preface}
\addcontentsline{toc}{chapter}{Preface}
\markboth{Preface}{Preface}

In November 1994 Leonard Adleman reported that DNA in a test tube had found a
route through a seven-city map of one-way roads that visits every city exactly
once.\cite{Adleman1994} The instance was tiny and the laboratory work took
about a week, but the idea was startling: a computation could be written
in the language of molecules, carried out by the chemistry of base pairing and
enzymes, and read back as evidence on a gel.

Three decades later the idea has grown into a field. Designed DNA strands now
implement logic circuits with hundreds of molecular species, chemical reaction
networks that behave like programs, neural networks that classify patterns and
even learn from molecular examples, self-assembling crystals that execute
algorithms as they grow, nanoscale shapes folded from a single scaffold strand,
archives that store digital files, and recorders that write the history of a
living cell into its own genome. In September 2026 a scaffolded DNA computer
whose correct answer is simply the most thermodynamically favoured assembly
was reported to execute programs such as binary arithmetic and graph
reachability.\cite{Sterin2026} This book explains how these systems work, why
they work, and what it takes to know that they have worked.

\section*{What is DNA computing?}
\Term{DNA computing} uses designed molecular interactions to carry out a
specified computation. Information is represented physically: by which
sequences are present, how many copies exist, which strands are bound to which
partners, and where molecules sit in space. Computation proceeds by physical
processes: complementary strands recognise each other, enzymes cut, join and
copy, strands displace one another, and tiles assemble. An answer is whatever
we can measure about the final molecular state and interpret against the
original question.

That definition already contains the three habits of mind this book asks
you to practise from the first page.
\begin{enumerate}
\item \textbf{Separate the specification, the mechanism and the evidence.}
A formal problem says what a correct answer is; molecules implement an
imperfect physical process; an instrument returns a signal that must be
interpreted. Most mistakes in molecular computing come from collapsing these
three levels into one.
\item \textbf{Count copies, not just identities.} A tube holds a
\emph{population}: many molecules, some identical, some missing, all finite.
Ideal set-theoretic reasoning becomes probability, concentration and noise.
\item \textbf{A model prediction is not a measurement.} We derive, simulate
and test models carefully, and we label every computed number as a model
output unless an experiment is cited.
\end{enumerate}

\section*{Why study it now?}
DNA computing sits where several revolutions meet. DNA synthesis and
sequencing have become cheap enough to write and read millions of designed
sequences. Thermodynamic models and design software predict hybridisation well
enough to engineer reliable interactions. Theoretical computer science has
supplied precise models (chemical reaction networks, tile assembly, molecular
automata) whose power and limits are now well understood. And biology itself
is the deepest example of molecular information processing: every cell
regulates, copies, repairs and records information encoded in nucleic acids.

The motivation is not to replace silicon. Electronic computers will remain
faster for almost every arithmetic task. Molecular computation matters where
electronics cannot go: inside a test tube of patient sample, on the surface of
a cell, within a nanostructure, or in an archive meant to last centuries. Its
most important applications are likely to be molecular decisions made in
molecular environments: diagnostics that compute on biomarkers, therapeutics
that act only when a logical condition holds, nanofabrication programmed by
sequence, and information storage at densities no tape can approach.

\section*{Who this book is for}
The book is written for undergraduates in biology, biochemistry,
bioengineering, chemistry, physics, computer science and electrical
engineering, and for graduate students and researchers entering molecular
programming. We assume secondary-school biology, algebra and some comfort with
reading a short program. Everything else (the chemistry of strands,
thermodynamics, kinetics, probability, automata, complexity and simulation) is
introduced where it is first needed. Two appendices provide compact primers:
\hyperref[app:biology]{Appendix~\ref*{app:biology}} for readers arriving from
computing, and \hyperref[app:math]{Appendix~\ref*{app:math}} for readers
arriving from the life sciences.

\section*{How the book teaches}
Every chapter is built around a small, concrete object that you can inspect
by hand: a six-vertex graph, a three-strand displacement reaction, a two-bit
sticker register, a five-tile crystal. The pattern is always the same.
\emph{Predict} what the object will do. \emph{Derive} the mechanism and its
mathematics with every intermediate step visible. \emph{Implement} a small
reference model. \emph{Check} it against an independent calculation.
\emph{Break} an assumption and see what fails. Only then do we generalise.

Illustrations are a principal teaching medium rather than decoration.
Molecular plates show strands, domains, enzymes and reaction intermediates in
successive states; algorithmic plates show graphs, state transitions and traces.
Wherever the two meet, a molecular panel is paired with its symbolic
representation so you can see what the abstraction keeps and what it discards.
The conventions are summarised in \hyperref[sec:conventions]{the visual
conventions} that follow this preface.

The implementation-first progression was informed by build-from-scratch
textbooks in other fields, notably Sebastian Raschka's \emph{Build a Large
Language Model (From Scratch)}, whose companion overview describes the
approach.\cite{RaschkaLLM} The prose, code, numerical examples, exercises and
illustrations of this book are original.

\section*{What is in the book}
Thirty-two chapters are grouped into eight parts.
\begin{description}
\item[Part I: Computing with molecules.] Adleman's experiment, the logic of
generate-filter-verify search, complexity, and the material cost of
molecular parallelism.
\item[Part II: The molecular toolkit.] DNA structure, hybridisation
thermodynamics, reaction kinetics, enzymes, PCR amplification, and
separation and measurement.
\item[Part III: Encoding and molecular algorithms.] Sequence design, filtering
algorithms, satisfiability constructions and sticker memories.
\item[Part IV: Formal models.] Splicing systems, automata, Watson--Crick
automata and universality.
\item[Part V: Programming molecular dynamics.] Toehold-mediated strand
displacement, chemical reaction networks as programs, digital, analog and
neural DNA circuits, and algorithmic self-assembly and DNA origami.
\item[Part VI: Reliability, scale and storage.] Noise and fault models,
resource limits, DNA data storage with error-correcting codes, and
reproducible laboratory evidence.
\item[Part VII: Modelling and simulation.] Symbolic, kinetic, stochastic and
differentiable models of molecular systems.
\item[Part VIII: From molecules to genomes.] Gene regulation and development,
modern molecular programming in realistic environments, and the bridge from
molecular to genomic computation.
\end{description}
An Epilogue surveys the research frontier and the open problems that a reader
of this book is now equipped to attack.

\section*{Companion code}
Each chapter has a small reference implementation in Python, with tests that
compare it against independent calculations. Printed listings are extracted
from those tested files, and numerical tables are generated from the same
results, so the book and its code cannot silently disagree. Readers who do not
program can skip the listings without losing the argument; readers who do
should run, modify and break them.

\section*{A note on scope and evidence}
This is an introductory scientific textbook. Its laboratory discussions
explain principles, controls and failure modes; they are not validated
protocols and do not replace supervised training. Executable models establish
results within their stated assumptions. Where the book reports experimental
findings it cites the primary literature, and it states when a source was read
only in part. The present text is a review edition: its technical content has
been checked by its authors and by executable tests, but independent
specialist review and classroom testing are still in progress. Corrections and
suggestions are warmly welcome.
